\section{Riesz estimates at fixed affine offsets}
\label{p:riesz-section}

An integrated density norm does not specify the value of a density on a
matrix-dependent graph of labels. We now construct a single density
version whose moment is controlled at every fixed affine offset. The
bound uses the supremum of ambient Riesz potentials of the original
measure. We first prove this analytic estimate, then use it on distance
shells to propagate finite-state continuation probabilities backward.
The shell argument retains an explicit logarithmic loss.

Throughout this section, fix an integer $d\ge16$ and put
\begin{equation}\label{p:riesz-parameters}
 k=q=\lfloor d/8\rfloor,\qquad \alpha=d/2.
\end{equation}
These choices give
\[
 q\ge2,\qquad d/16\le k=q\le d/8,\qquad
 \alpha-k-q\ge1,\qquad d-1-\alpha\ge1.
\]
The matrix $X$ will always denote a $k$ by $d$ standard Gaussian matrix.
Constants denoted by $C\ge1$ are absolute and may increase between uses.

\subsection{Ambient potentials and affine distances}

For a measurable $h:S^{d-1}\to[0,1]$, define the supremum of its
ambient Riesz potentials by
\begin{equation}\label{p:riesz-potential}
 V(h)=\sup_{z\in\R^d}
       \int_{S^{d-1}}\frac{h(s)}{\norm{s-z}^{\alpha}}\,d\sigma(s).
\end{equation}
For a measure $\mu$, the raw Riesz potential is
$U_\alpha^\mu(z)=\int\norm{s-z}^{-\alpha}\,d\mu(s)$
\cite[Section~1, p.~2]{CalefHardin2008v1}. Thus $V(h)$ is its supremum
over ambient centers for $\mu=h\sigma$; it is not the Riesz energy.
The kernel is understood as a nonnegative extended-valued function
at $s=z$. Write $B_r(z)=\{x\in\R^d:\norm{x-z}\le r\}$, and
intersect such balls with the sphere when integrating against
$\sigma$.

\begin{lemma}[Ambient potential bounds]\label{p:riesz-potential-finite}
For every $z\in\R^d$ and $0<r\le1/10$,
\begin{equation}\label{p:riesz-cap}
 \sigma(B_r(z))\le C^d r^{d-1}.
\end{equation}
Every measurable $h:S^{d-1}\to[0,1]$ has finite potential. Moreover, for
$0<\epsilon\le1/10$,
\begin{equation}\label{p:riesz-inner}
 \sup_{z\in\R^d}
 \int_{\norm{s-z}\le\epsilon}\norm{s-z}^{-\alpha}\,d\sigma(s)
 \le C^d\epsilon^{d-1-\alpha}.
\end{equation}
\end{lemma}

\begin{proof}
The ambient estimate \eqref{p:riesz-cap} follows from
Lemma~\ref{p:cap-mass}, with $2^{d-1}$ absorbed into $C^d$.
For dyadic outer radii $r_m=\epsilon2^{-m}$, the shell
$r_m/2<\norm{s-z}\le r_m$ contributes at most
$(2/r_m)^\alpha C^d r_m^{d-1}$. The point $s=z$ has zero
$\sigma$-measure, and
$\sum_{m\ge0}2^{-m(d-1-\alpha)}\le2$. Summing proves
\eqref{p:riesz-inner}. The same argument at radius $1/10$, together
with the kernel bound $10^\alpha$ outside that ball, proves
$V(h)<\infty$ uniformly over $0\le h\le1$.
\end{proof}

The potential bounds the measure of a tube around any affine subspace
through an ambient center. We use that fact to integrate successive
inverse distances when computing the Gaussian projection moment.

\Needspace{3\baselineskip}
\begin{lemma}[Inverse distance to an anchored affine space]
\label{p:riesz-distance}
Let $h:S^{d-1}\to[0,1]$ be measurable, $z\in\R^d$, $\rho>0$,
and let $F$ be an affine subspace through $z$ of dimension
$\ell\le q-1$.  Then $F\cap B_\rho(z)$ has zero $h\,d\sigma$
mass, and
\begin{equation}\label{p:riesz-distance-bound}
 \int_{B_\rho(z)}
       \operatorname{dist}(s,F)^{-k}h(s)\,d\sigma(s)
 \le C^d V(h)\rho^{\alpha-k}.
\end{equation}
\end{lemma}

\begin{proof}
For every ambient center $c$ and radius $r>0$, definition
\eqref{p:riesz-potential} gives
\begin{equation}\label{p:riesz-ball-control}
 \int_{B_r(c)}h\,d\sigma\le V(h)r^\alpha .
\end{equation}
Fix $0<r\le\rho$.  A maximal $r$-separated subset of
$F\cap B_\rho(z)$ has at most $(1+2\rho/r)^\ell
\le(3\rho/r)^\ell$ points, by comparing disjoint radius-$r/2$
balls in $F$ with the radius-$(\rho+r/2)$ ball.  It is an
$r$-net.  The orthogonal projection onto $F$ of any point in
$B_\rho(z)$ lies in $F\cap B_\rho(z)$, since $z\in F$.
Consequently the part of $B_\rho(z)$ within distance $r$ of $F$
is covered by the corresponding ambient radius-$2r$ balls.
Equation \eqref{p:riesz-ball-control} yields
\begin{equation}\label{p:riesz-tube}
 \int_{B_\rho(z)\cap\{\operatorname{dist}(s,F)\le r\}}
       h(s)\,d\sigma(s)
 \le 3^\ell2^\alpha V(h)\rho^\ell r^{\alpha-\ell}.
\end{equation}
This includes $\ell=0$, with a one-point net.  Since
$\alpha-\ell>0$, letting $r$ decrease to zero shows that
$F\cap B_\rho(z)$ has zero mass.

Every point of $B_\rho(z)$ is within distance $\rho$ of $F$.
Apply \eqref{p:riesz-tube} on the shells with outer distance
$r_m=\rho2^{-m}$ and bound the inverse power there by
$(2/r_m)^k$.  Their sum is at most
\[
 3^\ell2^{\alpha+k}V(h)\rho^{\alpha-k}
 \sum_{m\ge0}2^{-m(\alpha-\ell-k)}.
\]
Here
$\alpha-\ell-k\ge\alpha-(q-1)-k\ge2$,
so the series is at most $2$.  Since
$\ell\le d/8$, $\alpha=d/2$, and $k\le d/8$, the remaining
constant is at most $C^d$.  This proves
\eqref{p:riesz-distance-bound}.
\end{proof}

\subsection{A density at fixed affine offsets}

The next estimate uses $q$ points anchored at the chosen center $z$.
Their images under each Gaussian row have an anchored Gram matrix as
covariance. The inverse-distance bound just proved will integrate the
negative determinant power in the joint Gaussian density.

\begin{lemma}[Projection density at deterministic offsets]
\label{p:riesz-density}
Let $z\in\R^d$, $\rho>0$, let
$H\subset S^{d-1}\cap B_\rho(z)$ be measurable, and let
$h:S^{d-1}\to[0,1]$ be measurable. Fix these objects before drawing a
$k$ by $d$ standard Gaussian matrix $X$. The pushforward of
$h\mathbf 1_H\,d\sigma$ by $s\mapsto Xs$ admits a single jointly
Borel function
\[
 R:\R^{k\times d}\times\R^k\longrightarrow[0,\infty]
\]
such that $R(X,\cdot)$ is its Lebesgue density for every full-row-rank
$X$, and, for every fixed $b\in\R^k$,
\begin{equation}\label{p:riesz-density-bound}
 \left(\E_X R(X,Xz+b)^q\right)^{1/q}
 \le C^d V(h)\rho^{\alpha-k}.
\end{equation}
The same function $R$ is used for all $b$.  Moreover,
\begin{equation}\label{p:riesz-density-support}
 R(X,y)=0
 \quad\hbox{if}\quad
 \norm{y-Xz}>\norm X_{\mathrm{op}}\rho .
\end{equation}
\end{lemma}

\begin{proof}
Let $\mu$ be the finite Borel measure given by
$\mu(E)=\int_E h\mathbf1_H\,d\sigma$ on Borel subsets of the
sphere.  For $\eta>0$, let
$B_\eta^k=\{u\in\R^k:\norm u\le\eta\}$ and let
$\operatorname{vol}_k$ denote $k$-dimensional Lebesgue measure.  Define
\begin{equation}\label{p:riesz-ball-average}
 R_\eta(X,y)=
 \frac{1}{\operatorname{vol}_k(B_\eta^k)}
 \int \mathbf1_{\{\norm{Xt-y}\le\eta\}}\,d\mu(t).
\end{equation}
The integrand is Borel in $(X,y,t)$, so the parameter integral is
Borel in $(X,y)$.  We first prove
\eqref{p:riesz-density-bound} with $R_\eta$ in place of $R$,
uniformly in $\eta$.

Expand its $q$th moment at $Xz+b$ by Tonelli's theorem, using
$t_1,\ldots,t_q$ from the $q$ copies of $\mu$.  For a fixed tuple,
the columns $X(t_j-z)$ form an array whose $k$ rows are independent
centered Gaussian vectors with covariance
\[
 G=\bigl(\inner{t_i-z}{t_j-z}\bigr)_{1\le i,j\le q}.
\]
If $G$ is nonsingular, their joint density is bounded everywhere by
$(2\pi)^{-kq/2}(\det G)^{-k/2}$.  Therefore the probability that
all columns belong to the radius-$\eta$ ball about $b$ is at most
this bound times $\operatorname{vol}_k(B_\eta^k)^q$.

The singular tuples have zero $\mu^q$ mass.  Indeed, at the first
linearly dependent vector $t_j-z$, the point $t_j$ lies in
$z+\operatorname{span}(t_1-z,\ldots,t_{j-1}-z)$, an affine space
through $z$ of dimension at most $q-1$.  Lemma~\ref{p:riesz-distance}
assigns zero $\mu$ mass to it.  The degeneracy set is measurable,
and Fubini's theorem applied at each $j$ proves the claim.  For a
nonsingular tuple, Gram--Schmidt triangularizes the matrix with
columns $t_j-z$ with diagonal entries equal to the successive
orthogonal distances.  Taking the determinant of its Gram matrix
gives
\begin{equation}\label{p:riesz-gram}
 (\det G)^{-k/2}
 =\prod_{j=1}^q
 \operatorname{dist}\!\left(
 t_j,z+\operatorname{span}(t_1-z,\ldots,t_{j-1}-z)
 \right)^{-k}.
\end{equation}
Integrate this product first in $t_q$, then in $t_{q-1}$, and so on.
For every nonsingular prefix, Lemma~\ref{p:riesz-distance} bounds
the next integral by $C^dV(h)\rho^{\alpha-k}$.  The exceptional
prefixes have zero product measure by the same degeneracy argument.
It follows that the integral of
\eqref{p:riesz-gram} is at most
$(C^dV(h)\rho^{\alpha-k})^q$.  The ball volumes cancel in the
moment expansion, and the Gaussian prefactor is at most one.
Thus
\[
 \E_X R_\eta(X,Xz+b)^q
 \le \bigl(C^dV(h)\rho^{\alpha-k}\bigr)^q
\]
for every $\eta>0$ and every fixed $b$.

Now choose the function once and for all:
\begin{equation}\label{p:riesz-density-version}
 R(X,y)=\liminf_{m\to\infty}R_{1/m}(X,y).
\end{equation}
It is jointly Borel. To see that it is a density, fix a full-row-rank
$X$. Set $D=(XX^{\mathsf T})^{1/2}$ and $P=D^{-1}X$, so $D$ is invertible,
$P$ has orthonormal rows, and $X=DP$.
Lemma~\ref{p:coordinate-density} makes $P_\#\sigma$ absolutely
continuous, and the invertible map $D$ preserves absolute continuity.
Since $\mu\le\sigma$, the pushforward $X_\#\mu$ is also absolutely
continuous. The Lebesgue differentiation
theorem~\cite[Theorem~1.6.19]{TaoMeasure2011}, applied to its
$L^1$ density, shows that
\eqref{p:riesz-density-version} equals that density almost
everywhere in $y$.  A Gaussian $X$ has full row rank almost
surely, since each successive Gaussian row avoids the span of
the preceding rows.

For each fixed $b$, the identity
$R(X,Xz+b)=\liminf_m R_{1/m}(X,Xz+b)$ holds by the definition of
the single function $R$.  Fatou's
lemma~\cite[Corollary~1.4.47]{TaoMeasure2011} over $X$ now gives
\eqref{p:riesz-density-bound}.  These are separate moment
statements for each deterministic $b$; no simultaneous
almost-sure bound in $X$ over all offsets is asserted.  Finally,
if $\norm{y-Xz}>\norm X_{\mathrm{op}}\rho$, then every $Xt$ with
$t\in H$ is at distance at least
$\norm{y-Xz}-\norm X_{\mathrm{op}}\rho>0$ from $y$.  The averages in
\eqref{p:riesz-ball-average} are eventually zero, proving
\eqref{p:riesz-density-support}.
\end{proof}

The explicit choice \eqref{p:riesz-density-version} is the feature
needed for the next step. A density chosen only up to a
Lebesgue-null set in $y$ for each $X$ would not justify the moment
estimate at the matrix-dependent point $Xz+b$.
An integrated $L^q(dX\,dy)$ estimate controls a density only up to
product-null sets, while the graph $y=Xz+b$ has zero product measure.
Positive cap domination instead controls the measure of every Borel
set by a positive sum of cap measures.  Neither conclusion supplies
the specified-offset moment in Lemma~\ref{p:riesz-density}; its
anchored determinant estimate and explicit limit are therefore kept
as a separate argument.

\subsection{Continuation probabilities and the terminal potential}

We apply the specified-offset estimate to the finite-state experiment of
Section~\ref{p:geometry-model}. Fix nonnegative integers $M,T$ and an
accuracy $0<\epsilon\le1/10$, and write
\[
 p_{\mathrm{succ}}=\Prob\{\arccos\inner{\widehat S}{S}\le\epsilon\},
 \qquad L=\log(1/\epsilon),\qquad
 N=\lceil T/k\rceil,\qquad W=(T+2)2^M.
\]
The success probability averages over the uniform signal, Gaussian
features, and learner randomness. The factor $W$ will count the retained
states together with their possible terminal stopping indices.

\begin{proposition}[Riesz success estimate]\label{p:riesz-success}
There is an absolute constant $C_0\ge1$ such that every learner just
specified satisfies
\begin{equation}\label{p:riesz-success-bound}
 p_{\mathrm{succ}}\le B J^N,\qquad
 B=C_0^d\epsilon^{\,d/2-1},\qquad
 J=1+C_0^d(1+L)W^{1/q}.
\end{equation}
This includes $T=0$.  No asymptotic assumption on $M$ is needed for
this estimate.
\end{proposition}

We first bound the potential of each terminal output law. To control its
increase across a block, we apply Lemma~\ref{p:riesz-density} at $Xz+b$
and then integrate over $b$. The factor $1+L$ in
\eqref{p:riesz-success-bound} counts the distance shells and remains
in the final parameter argument.

Use the fixed-seed and Borel-kernel convention of
Section~\ref{p:geometry-model}: first fix an admissible shared seed,
under which the initial state is independent of the signal and rows,
and then apply Lemma~\ref{p:borel-kernels} to choose Borel rules. These
rules define continuations from every state. All bounds below are uniform
in the fixed rules; after proving the conditional bound, we average the
original conditional success probabilities over the seed.

Complete unused rules and terminal laws arbitrarily, and force a halt
by $T$.  After a halt, retain a tag consisting of the stopping index
and terminal state and ignore further samples.  At any time there are
at most $2^M$ live states and at most $(T+1)2^M$ such tags, including
a possible immediate halt. Hence the stated bound $W$ applies. Add
independent unused Gaussian rows after $T$ to complete the last block of length $k$.  Let
$\mathcal V_i$ be the state set at the boundary after block $i$,
$0\le i\le N$; thus $\lvert\mathcal V_i\rvert\le W$.

For every $v\in\mathcal V_i$ and $s\in S^{d-1}$, let $h_i(v,s)$ be
the success probability of the remaining execution started at that
boundary state.  This definition also applies to states not reached
by the original run.  At $i=N$ it is the success probability of the
fixed output law associated with the terminal state.  These functions
are Borel and take values in $[0,1]$: the terminal assertion follows
by integrating the Borel success indicator against the output law,
and the assertion at earlier boundaries follows inductively from
the formula below.  If $X$ is a fresh $k$ by $d$ standard Gaussian
matrix, then for $1\le i\le N$
\begin{equation}\label{p:riesz-recursion}
 h_{i-1}(v,s)
 =\E_X\sum_{w\in\mathcal V_i}K^{(i)}_{v,w}(X,Xs)h_i(w,s).
\end{equation}
Here $K^{(i)}_{v,w}(X,y)$ is the probability of the block transition
from $v$ to $w$ on the exact row pairs $(X,y)$.  These kernels are
Borel, nonnegative, and sum to one in $w$ for every $(X,y)$.
They incorporate internal randomization and the ignored rows after a
halt.  The functions $h_i(w,\cdot)$ are fixed independently of $X$.
The fresh Gaussian block is independent of the signal, initialization,
past rows, and past coins. Starting the continuation at the specified
state therefore gives exactly \eqref{p:riesz-recursion}; the state and
boundary index contain all retained sample-dependent information.

\begin{lemma}[Terminal potential]\label{p:riesz-terminal-lemma}
Let $Q$ be a signal-independent probability law on unit-vector outputs,
and let $h(s)=Q\{u:\arccos\inner{u}{s}\le\epsilon\}$. Then
\begin{equation}\label{p:riesz-terminal}
 V(h)\le C^d\epsilon^{d-1-\alpha}.
\end{equation}
\end{lemma}

\begin{proof}
Lemma~\ref{p:cap-mass} gives $\int h\,d\sigma\le\epsilon^{d-1}$,
also for a randomized output law. For any ambient center $z$, use
$h\le1$ and \eqref{p:riesz-inner} inside $B_\epsilon(z)$. Outside
that ball, bound the kernel by $\epsilon^{-\alpha}$ and use the total
mass bound. The sum is at most $C^d\epsilon^{d-1-\alpha}$ uniformly
in $z$, proving \eqref{p:riesz-terminal}.
\end{proof}

Set
\[
 V_i=\max_{v\in\mathcal V_i}V(h_i(v,\cdot)).
\]
The maximum is finite by Lemma~\ref{p:riesz-potential-finite}.
Lemma~\ref{p:riesz-terminal-lemma} gives $V_N\le C^d\epsilon^{d-1-\alpha}$.
We next estimate the contribution of a single distance shell in
\eqref{p:riesz-recursion}.

\subsection{One block and the distance shells}

\begin{lemma}[One shell across a block]\label{p:riesz-shell}
For $1\le i\le N$, $v\in\mathcal V_{i-1}$, $z\in\R^d$, and
$\rho>0$, let
$H_\rho(z)=\{s\in S^{d-1}:\rho/2<\norm{s-z}\le\rho\}$.
Then
\begin{equation}\label{p:riesz-shell-bound}
 \int_{H_\rho(z)}
   \frac{h_{i-1}(v,s)}{\norm{s-z}^{\alpha}}\,d\sigma(s)
 \le C^d W^{1/q}V_i.
\end{equation}
\end{lemma}

\begin{proof}
For each successor $w$, use the particular density $R_w(X,y)$
of Lemma~\ref{p:riesz-density} for the measure
$h_i(w,\cdot)\mathbf1_{H_\rho(z)}\,d\sigma$.
Tonelli's theorem, \eqref{p:riesz-recursion}, and pushforward
integration give
\begin{equation}\label{p:riesz-shell-pushforward}
\begin{split}
 \int_{H_\rho(z)}
    \frac{h_{i-1}(v,s)}{\norm{s-z}^{\alpha}}\,d\sigma(s)
 &\le (2/\rho)^\alpha
       \E_X\int_{\R^k}
       \sum_wK^{(i)}_{v,w}(X,y)R_w(X,y)\,dy\\
 &\le (2/\rho)^\alpha
       \E_X\int_{\norm{y-Xz}\le\norm X_{\mathrm{op}}\rho}
       \max_wR_w(X,y)\,dy .
\end{split}
\end{equation}
The pushforward identity used in the first inequality holds for every
full-row-rank $X$, hence almost surely.  The second inequality uses
the probability-vector property of $K^{(i)}$ and
\eqref{p:riesz-density-support}.

We need a deterministic integration region after averaging over
$X$.  Choose a sufficiently large absolute $D_0$ and partition
the matrices into
\[
 E_1=\{\norm X_{\mathrm{op}}\le D_0\sqrt d\},\qquad
 E_j=\{D_0(j-1)\sqrt d<\norm X_{\mathrm{op}}\le D_0j\sqrt d\}
 \quad(j\ge2).
\]
These bins obey
\begin{equation}\label{p:riesz-bins}
 \P(E_j)\le \exp\bigl(-4(j-1)^2d\bigr)\qquad(j\ge2).
\end{equation}
For completeness, take $1/4$-nets on the unit spheres of $\R^d$
and $\R^k$.  A maximal separated set has at most $9^d$ and
$9^k$ points, respectively, by the disjoint-ball volume argument.
Approximating each unit vector by its net point incurs an error
at most $\norm X_{\mathrm{op}}/2$ in a bilinear form.  Hence
$\norm X_{\mathrm{op}}$ is at most twice the largest absolute
bilinear form on the two nets.  Each such form is standard
Gaussian, so its elementary tail bound and a union bound give
\begin{equation}\label{p:riesz-operator-tail}
 \P\{\norm X_{\mathrm{op}}>t\}
 \le 2\cdot9^{d+k}\exp(-t^2/8).
\end{equation}
Putting $t=D_0(j-1)\sqrt d$ proves
\eqref{p:riesz-bins} after increasing $D_0$, since $k\le d/8$.

For every deterministic $b\in\R^k$, H\"older's inequality and
\eqref{p:riesz-density-bound} give
\begin{equation}\label{p:riesz-bin-holder}
\begin{split}
 \E_X\!\left[\mathbf1_{E_j}\max_wR_w(X,Xz+b)\right]
 &\le \P(E_j)^{1-1/q}
       \left(\sum_w\E_X R_w(X,Xz+b)^q\right)^{1/q}\\
 &\le \P(E_j)^{1-1/q}W^{1/q}
       C^dV_i\rho^{\alpha-k}.
\end{split}
\end{equation}
No conditional version of the density lemma is being used:
H\"older's inequality pays for the bin event.  For a fixed matrix
in $E_j$, translate $y=Xz+b$ in
\eqref{p:riesz-shell-pushforward}; the remaining region is
$\norm b\le D_0j\sqrt d\,\rho$.  Joint measurability permits
Tonelli's theorem to integrate \eqref{p:riesz-bin-holder} over
this deterministic ball.

The volume of that ball is at most $C^d j^k\rho^k$.  Indeed, the unit-ball bound \eqref{p:unit-ball-volume} gives
$v_k\le(2\pi e/k)^{k/2}$. Since $d/k\le16$,
$v_k(D_0j\sqrt d\,\rho)^k\le C^d j^k\rho^k$.
Also, by $q\ge2$ and \eqref{p:riesz-bins},
\[
 \sum_{j\ge1}j^k\P(E_j)^{1-1/q}
 \le 1+\sum_{j\ge2}
       \exp\bigl(k\log j-2(j-1)^2d\bigr)
 \le 2 .
\]
For the last bound use
$k\log j\le(d/8)(j-1)^2$ for $j\ge2$ and sum the resulting
geometric upper bound.  Combining the volume and bin estimates
with \eqref{p:riesz-shell-pushforward} proves
\eqref{p:riesz-shell-bound}.  The powers of the shell radius
cancel exactly:
\[
 \rho^{-\alpha}\rho^{\alpha-k}\rho^k=1.
\]
\end{proof}

\begin{proof}[Proof of Proposition~\ref{p:riesz-success}]
Outside $B_\epsilon(z)$ use the shells
\[
 2^m\epsilon<\norm{s-z}\le2^{m+1}\epsilon,
 \qquad m=0,1,\ldots .
\]
For every ambient $z$, at most $C(1+L)$ of these meet the
sphere.  If $\norm z\le3$, all distances are at most $4$, giving
at most $\lceil\log_2(4/\epsilon)\rceil$ nonempty shells.  If
$\norm z>3$, the distances lie between $\norm z-1$ and
$\norm z+1$, whose ratio is less than $2$; only a bounded number
of dyadic shells can then meet that interval.

Inside $B_\epsilon(z)$ use $h_{i-1}\le1$ and
\eqref{p:riesz-inner}.  On each remaining shell use
Lemma~\ref{p:riesz-shell}.  Taking the supremum over $z$ and the
maximum over $v$ gives
\begin{equation}\label{p:riesz-step}
 V_{i-1}\le C^d\epsilon^{d-1-\alpha}
          +C^d(1+L)W^{1/q}V_i .
\end{equation}
Choose $C_0$ large enough that the first term here and the
terminal value $V_N$ are at most $B$ in
\eqref{p:riesz-success-bound}, and that the second coefficient is
at most $J-1$.  With $U_i=\max\{V_i,B\}$, equation
\eqref{p:riesz-step} gives $U_{i-1}\le J U_i$, while $U_N=B$.
Hence $V_0\le BJ^N$, also when $N=0$.

At $z=0$ the kernel in \eqref{p:riesz-potential} equals one on
the unit sphere, so $\int h_0(v,s)\,d\sigma(s)\le V_0$ for
each initial state. For almost every fixed shared seed, averaging over its
conditional initial-state law preserves the bound, because that law is
independent of the signal and rows. Averaging the original conditional
success probabilities over the seed then
preserves it again.  This proves \eqref{p:riesz-success-bound}.
\end{proof}

\subsection{The logarithmic lower bound}

We now prove Theorem~\ref{p:regression} from the Riesz success estimate.
The parameter argument retains the stopping-index factor in $W$ and
the shell-count factor $1+L$. Terminal-tag counting controls both in
the only case still requiring proof, namely $T<dL$.

\begin{proof}[Proof of Theorem~\ref{p:regression} by the Riesz method]
There is a useful bound that counts possible terminal tags.
For fixed rules as above, write $p_{\mathrm{succ},\mathrm{rule}}$ for
the conditional success probability.
For a tag $a$, let $f_a(s)$ be the probability of reaching it
given $s$, and let $g_a(s)$ be its output law's success
probability.  This output law depends only on the tag, and
$0\le f_a\le1$.
Lemma~\ref{p:cap-mass} gives
$\int g_a\,d\sigma\le C^d\epsilon^{d-1}$, including a randomized
output law.  Sum over at most $W$ tags to obtain
\begin{equation}\label{p:riesz-packing}
 p_{\mathrm{succ},\mathrm{rule}}=\sum_a\int f_a(s)g_a(s)\,d\sigma(s)
 \le W C^d\epsilon^{d-1}.
\end{equation}
The right side is independent of the chosen rules, so averaging
also gives $p_{\mathrm{succ}}\le W C^d\epsilon^{d-1}$.

Write $L=\log(1/\epsilon)$.  If $T\ge dL$, the conclusion
already holds.  In the remaining case $T<dL$, assume $p_{\mathrm{succ}}\ge2/3$.
Taking logarithms in the averaged form of \eqref{p:riesz-packing} gives
\[
 (d-1)L
 \le M\log2+\log(T+2)+d\log C+\log(3/2).
\]
For the given memory sequence, eventually $M\le d^2$.  Since
$L\ge\log10$ and $T<dL$, we also have
\[
 \log(T+2)\le\log(2d)+\log L\le\log(2d)+L
\]
for large $d$.  It follows that
$(d-2)L\le d^2\log2+d\log C+\log(2d)+\log(3/2)$.
Consequently, only in this remaining case,
\begin{equation}\label{p:riesz-short-run}
 L\le Cd,\qquad \log W=M\log2+\log(T+2)\le Cd^2 .
\end{equation}
Thus the multiplier from Proposition~\ref{p:riesz-success} obeys
\[
 \log J
 \le \log2+d\log C_0+\log(1+L)+\frac{\log W}{q}
 \le Cd ,
\]
using $q\ge d/16$.  In particular the shell-count term has
been bounded rather than removed.

Now \eqref{p:riesz-success-bound} and $p_{\mathrm{succ}}\ge2/3$ imply
\begin{equation}\label{p:riesz-final}
 (d/2-1)L
 \le d\log C_0+N\log J+\log(3/2)
 \le Cd(1+N)+\log(3/2)
 \le C'(d+T),
\end{equation}
because $N=\lceil T/k\rceil\le16T/d+1$.
The common full-data bound, Lemma~\ref{p:haar-linear}, gives
$T>d/8$; hence the additive $d$ on the right of
\eqref{p:riesz-final} is at most $8T$.
For $d\ge4$ the left coefficient is at least $d/4$, so
$T\ge c\,dL$ for an absolute $c>0$. This combines with the
case $T\ge dL$ already separated.
\end{proof}
